\documentclass[10pt]{article}
\usepackage[margin=20mm]{geometry}
\usepackage{amsmath,amssymb,fontspec,unicode-math,xurl,hyperref}
\IfFontExistsTF{DejaVu Serif}{\setmainfont{DejaVu Serif}}{\setmainfont{Latin Modern Roman}}
\setmathfont{latinmodern-math.otf}
\setlength{\parindent}{0pt}\setlength{\parskip}{6pt}\sloppy
\begin{document}
Scientific verification report: 0044\par
Position: Retrieval-augmented systems can be dangerous medical communicators\par
I25-40149 | PMLR 267 (ICML 2025) | AI-assisted review | 24 October 2026\par
Supported outcome: Minor: a qualitative mastectomy example misstates its released context, and three source-overlap rows conflict with their raw citations. Tables3-4 and the published Jaccard summary means are reproducible from released aggregates; broad output-feature trends survive the bounded corrections.\par
\# 1. Identity, reviewed version and scope\par
Lionel Wong, Ayman Ali, Raymond Xiong, Shannon Shen, Yoon Kim and Monica Agrawal. PMLR267:82347-82359 (ICML2025). The publisher metadata renders Shen’s given name as Zejiang; the original PDF uses Shannon. Identity is fixed by the exact publisher-linked PDF, title and full author group, not this metadata name variant.\par
Reviewed papers/0044/0044.pdf: 13 pages, 835,933 bytes, SHA-256 137f64c7a9d5b6e4e2af07f69bc58ceafb1537264cf83d622dd6413800d5d3c0, matching the controller receipt. Page references mean printed pages 1-13, equal to one-based PDF pages. The official arXiv record lists 2502.14898v1 (18 February2025,12pages) and v2 (9 June2025,13pages); v2 is the latest listed version. Both were archived separately; the publisher original was preserved.\par
Read Sections1-6, Figures1-2, Tables1-5 and AppendicesA-C. Recomputed published summary statistics from released files, inspected the annotation source without executing it, and checked a selected medical communication example against its complete released answer and primary clinical context. No exhaustive medical validation of all generated answers or all 13 disputed diagnoses was attempted.\par
The official publisher Software link and paper point to rayarxti/rag-medical-communicator. The meaningful 25-file tree was preserved at commit74c37eb584d5c8f8d0d99ca4f0c63714c2ae820f (14February2025), with Git blob verification and SHA-256 acquisition manifests. All code/data analysis was static or reviewer-written. Human scientific decision and author adoption remain pending.\par
\# 2. Quantitative results reproduced and their denominators\par
The paper uses selected query templates for disputed diagnoses and procedure safety/complications. The reported measures are properties of captured answers and cited-source sets. They are not population estimates of patient harm. The following values are reproducible from released per-topic/template proportion files and the released Jaccard table:\par
Reported result | Google AIO: reconstructed | Perplexity: reconstructed\par
Table3, controversy mentioned for symptom queries:56\% /69\% | Mean14/25 nonmissing cells=56.00\% | 18/26cells=69.23\%\par
Table4, complication statistics:4\% /5\% | Mean82cells=3.8618\% | 4/83cells=4.8193\%\par
Table4, procedure benefits:6\% /10\% | Mean82cells=6.3415\% | 8/83cells=9.6386\%\par
Mean source Jaccard:0.16 /0.31 | 84topic-days:0.1635779 | 28topic-days:0.3145550\par
Direct-condition queries also retain the reported100\% controversy-mention result in the inspected labels. For Google, the summary cells can themselves represent averages across dates. The table reconstruction uses equal-weight nonmissing topic/template cells; it must not be described as simple pooling of all captured answer rows. The exact code and counts are preserved in independent\_quantitative\_checks.py and independent\_quantitative\_results.json.\par
Pooling the released complication-answer rows instead produces Google statistics9/225=4.00\% and benefits16/225=7.11\%, and Perplexity statistics4/89=4.49\% and benefits8/89=8.99\%. These different quantities do not refute the reproduced table means. They show why the unit of analysis, weighting, missingness and raw-to-summary mapping need to be explicit.\par
The raw category exports include duplicates and some mismatched query strings, including a duplicated direct Wind turbine syndrome query and duplicated procedure queries. The combined analyzed exports contain177 distinct query strings rather than directly recovering the stated178 intended queries; the complication export has89 rows per single collection rather than83 unique intended strings. A full collection and aggregation script is absent. Treat exact query lineage and the83\% availability statement as incompletely reconstructed, not as proof that the published labels or substantive trends are false.\par
Scope clarification for source overlap. Google has28 of84 topic-day pairs with at least one recorded source set empty:9 have both sets empty and19 have only one empty. The table assigns0.0 when both are empty; that requires an explicit convention because the usual intersection/union fraction is0/0. Conditioning on both recorded source sets being nonempty gives0.2453668 over56pairs. Perplexity has no such empty-set pairs. This is a sensitivity analysis for a different conditional estimand, not an automatic replacement of the published mean. Empty source sets are not uniformly synonymous with absent answers; three uniquely matched answers actually contain numbered citation URLs despite zero source counts in the summary (Problem2).\par
The reproduced differences support the narrower claim that query wording changes answer features and cited material in this selected dataset. They do not by themselves quantify how dangerous an answer is or whether a different citation set is inappropriate.\par
\# 3. Problem1: the double-mastectomy case omits context from its own evidence\par
Source claim and location. Section3.2,p6, uses the query about whether to get double mastectomy for cancer as an example of missing patient-selection conditions. It excerpts a personal-choice statement, then contrasts it with the original source’s purported absence of preventive benefit except for high-risk patients with specific mutations. Section4,p9, revisits the emotional implications of that phrase.\par
Decisive response check. The repository file google-ai-overview-misc/should-i-get-double-mastectomy-for-cancer, SHA-256 30ff8323b16ac5a0318e8dc8d15735493673c18172eaf52d35af6a3c80514293, preserves the date string12-06-2024 without requiring a locale interpretation. Before the personal-choice bullet, its full answer states that suitability depends on individual risk, advises against double mastectomy for cancer in one breast without genetic predisposition, offers less extensive alternatives, and distinguishes strong-family-history/BRCA situations. It also notes incomplete risk elimination and later discusses adverse effects.\par
Thus the complete released answer contains the very selection caveat portrayed as absent. The isolated personal-choice sentence could still be criticized for wording, salience or emotional impact. That is a narrower claim than omission of the condition from this answer. This audit does not certify every statement in the answer as medically correct; other inaccurate descriptions would not establish the specified omission.\par
Clinical endpoint check. Preventing a new cancer in the opposite breast, improving survival and achieving a favorable overall benefit-harm balance are different endpoints. The answer’s NCI source explicitly distinguishes contralateral risk reduction from lack of demonstrated mortality reduction. Northwestern Medicine’s cited clinician discussion focuses on survival, and Safer Care Victoria distinguishes low baseline contralateral risk from life-expectancy benefit. The current pages are not verified snapshots from the original collection date.\par
Contemporaneous primary evidence also supports the distinction: Giannakeas, Lim and Narod (JAMA Oncology2024; DOI10.1001/jamaoncol.2024.2212) report greatly reduced contralateral cancer incidence with bilateral mastectomy but similar breast-cancer mortality in matched cohorts. This observational evidence does not establish a universal causal treatment recommendation. It does show why absence of an established survival benefit cannot be restated as absence of every preventive effect.\par
Alternative reading and impact. The authors may intend preventive benefit to mean clinically worthwhile net benefit, and may intend to criticize the isolated phrase rather than the answer as a whole. Those readings should be stated explicitly. As written, the full-response omission claim is unsupported by the released evidence, and the endpoint language is too broad. Minor is warranted for this bounded qualitative case; the correction does not overturn the other cases, the reconstructed summary results, or the general possibility of factually accurate but misleading communication.\par
Scoped repair. Show the full answer or accurately state its preceding patient-selection caveats. If retaining the case, focus the criticism on the precise residual implication and distinguish risk reduction, survival and net benefit. No patient-belief experiment or clinical outcome reanalysis was conducted. A different unpublished response could change the case assessment, but would require its exact timestamped artifact. Evidence: pragmatics\_medical\_memo.json, the preserved response, and medical\_source\_validation.json.\par
\# 4. Problem2: three released source sets are incorrectly recorded as empty\par
Source claim and location. Section2.2,p4 reports mean source Jaccard0.16 for Google. The released results/safe-dangerous/jaccard\_by\_date.csv reproduces0.1635779. However, three uniquely matched raw safe/dangerous response pairs disagree with their summary source counts.\par
Complete static check. For each case, match method, topic, query wording and date in output\_csv\_safe\_dangerous.csv; exactly one raw row exists for each side. Convert literal newline escapes to newlines and extract unique exact URLs from numbered citation entries. No URL canonicalization, web refetch or author code is involved. The source-set definition is explicit and both sides contain citations.\par
Case | Released safe count / J | Raw safe / dangerous counts | Intersection / union | Recomputed J\par
Breast implants,20Dec2024 | 0 / 0 | 13 / 6 | 3 / 16 | 0.187500\par
Breast implants,21Dec2024 | 0 / 0 | 11 / 5 | 3 / 13 | 0.230769\par
Liver transplant,20Dec2024 | 0 / 0 | 8 / 5 | 3 / 10 | 0.300000\par
These are not simply no-answer or no-citation cases: the raw answers contain13,11 and8 distinct numbered source URLs. Exact intersection members are preserved in independent\_jaccard\_three\_cases.json. For example, the first pair shares two distinct FDA URLs and one UCLA Health URL, yielding3/(13+6-3)=3/16.\par
Changing only these three zeros, while retaining all other author rows and conventions, moves the84-pair mean from0.1635779 to0.1721287, which rounds to0.17. This verifies a limited arithmetic repair, not a complete reconstruction of every extracted source set or an estimate under a new collection protocol.\par
Alternative reading and provenance. The release has no source-extraction or aggregation code, so an undisclosed filtering rule or a different historical input cannot be ruled out. Those possibilities require evidence; the current raw/summary pair is internally inconsistent under the explicit ordinary URL-set calculation. The audit does not infer that the available script generated the published table, nor invent a cause for the discrepancy.\par
Impact and judgment. Minor, a bounded released-data/source-overlap discrepancy. The local repair changes the rounded Google mean but preserves relatively low overlap and its ordering against Perplexity0.3145550. No material overturning of the broad communication-risk thesis is demonstrated. Preserve the originals, supply the actual extraction lineage, and correct affected summaries if these raw records are the intended inputs.\par
\# 5. Evidence that survives and limits of interpretation\par
The core logical distinction is sound: individual accurate propositions can produce a misleading impression through selection, ordering or omitted context. The selected safe/dangerous examples and aggregate answer-feature changes give reasons to investigate that risk. Two differently framed questions can also legitimately call for different facts; low source overlap alone does not establish communication failure.\par
The paper explicitly acknowledges on pp4-5 that actual effects on users’ beliefs need further study. Its quantitative labels detect controversy mentions, numerical frequency information and benefits, rather than measuring patient interpretation or downstream harm. Accordingly, this review does not invent a requirement for a new patient experiment or treat its absence as a demonstrated error. General claims about population prevalence, superiority of conventional search or causal patient harm should remain limited to the evidence actually collected.\par
Some qualitative examples contain literal factual errors, such as the rope-worm/pinworm conflation or anatomical confusion. Those examples cannot by themselves isolate the narrower phenomenon of accurate-but-misleading output. The paper, however, also acknowledges egregious misinterpretations and discusses a wider set of communicative failures. Their inclusion is not a disproof of the central possibility claim.\par
Sections4-5 propose pragmatic modeling, dialogue and interface approaches rather than reporting validated solutions. The alternative-view discussion acknowledges that some users may already know the benefits and want a focused answer. This weakens any universal inference from a missing benefit sentence to harm, but is already part of the paper’s stated debate. No new experimental requirement is imposed on these normative recommendations.\par
\# 6. Annotation provenance and version history\par
Released evaluator.py contains three GPT-4o judgment functions, with an overrideable model parameter and defaultgpt-4o-2024-11-20. These inspect controversy, frequency statistics and benefits. Static inspection found no evidence that this source was executed during the audit; it was not imported or run.\par
The historicalv1 specifies that GPT-4o version and spot checks. Thev2/publisher methods add two GPT-4o versions and manual adjudication of disagreements. The public repository’s last commit predates both paper versions. Its CSVs expose final labels but not separate judge outputs, the second model identifier or an adjudication log; the tree also lacks collection/orchestration/aggregation scripts.\par
This is a Needs Evidence provenance limitation for independently reconstructing the final judging pipeline. It is not evidence that a second judge or human adjudication never occurred: the model argument can be overridden, and the release may predate later work. The headline rates persist between the inspected versions, but that alone does not prove identical annotation histories.\par
The official version record contains no later arXiv revision. No correction was found in the bounded official pages/repository checks. Complete OpenReview discussion/revision material was not accessible. These are limited search outcomes, not an assurance that no correction exists anywhere. The mastectomy wording persists in v1 printedp5, v2p6 and the publisherp6.\par
\# 7. Conclusion and correction status\par
Item | Supported judgment / verification | Remaining work\par
Double-mastectomy example | Minor; complete response and endpoint distinction checked. | Revise case framing; author adoption and human decision pending.\par
Three Jaccard source-set rows | Minor; exact URL intersections verified; limited mean repair0.16358 to0.17213. | Full extraction lineage and all-row validation remain open; original data preserved.\par
Tables3-4 and Jaccard summary means | Aggregate means reproduced; this does not validate all raw-to-summary transformations. | Clarify weighting, missing-set convention and duplicate handling.\par
Judging pipeline and query mapping | Needs Evidence; missing provenance does not raise severity. | Provide second judge ID/labels, adjudication record and collection/aggregation mapping.\par
Patient effects and proposed remedies | Explicitly unmeasured / proposed; scope clarification. | Fresh appropriately designed research if these claims are to be quantified.\par
The supported overall severity is Minor. Two bounded defects are established: one qualitative illustration and three released overlap rows. Substantial quantitative evidence survives their correction. No Major has been demonstrated within the inspected scope. The report does not declare whole-paper correctness, validate every clinical example, or provide patient-specific medical advice.\par
\# 8. Primary references and evidence identifiers\par
Wong etal.(2025), PMLR267:82347-82359; arXiv2502.14898v1 andv2.\par
https://proceedings.mlr.press/v267/wong25b.html\par
https://arxiv.org/abs/2502.14898\par
Author release, commit74c37eb584d5c8f8d0d99ca4f0c63714c2ae820f; preserved meaningful tree, not executed.\par
https://github.com/rayarxti/rag-medical-communicator/tree/74c37eb584d5c8f8d0d99ca4f0c63714c2ae820f\par
NCI, Surgery to Reduce the Risk of Breast Cancer, current retrieved source [1] in the preserved generated answer; historical snapshot not established.\par
https://www.cancer.gov/types/breast/risk-reducing-surgery-fact-sheet\par
Northwestern Medicine, Do You Need a Double Mastectomy?, source[3]; current page displays Updated July2024, which does not prove unchanged historical text.\par
https://www.nm.org/healthbeat/healthy-tips/deciding-to-get-double-mastectomy\par
Safer Care Victoria, Double mastectomy for cancer in a single breast, source[5]; current official page read through web retrieval.\par
https://www.safercare.vic.gov.au/best-practice-improvement/clinical-guidance/non-urgent-elective-surgery/double-mastectomy-for-cancer-in-a-single-breast\par
Giannakeas,V.; Lim,D.W.; Narod,S.A.(2024). Bilateral Mastectomy and Breast Cancer Mortality. JAMA Oncology10(9):1228-1236. DOI10.1001/jamaoncol.2024.2212; PMID39052262. Primary abstract inspected; no cohort-data reanalysis.\par
https://pubmed.ncbi.nlm.nih.gov/39052262/\par
\# Appendix A. Actual work, failures and provenance\par
Identity/hash verification preceded analysis. Reviewer scripts independent\_csv\_structure.py and independent\_quantitative\_checks.py read the preserved CSVs using ordinary data parsing and arithmetic. They did not import author evaluator.py or call a model. Outputs record row counts, duplicates, nonmissing-cell means, response-pooled alternatives and source-set sensitivity calculations. Source pp4 and6 were rendered with preinstalled Poppler and visually inspected. independent\_jaccard\_three\_cases.py separately checks the three uniquely matched source lists and the limited mean correction.\par
Medical-source acquisition succeeded for Northwestern and NCI HTML. The local Safer Care Victoria download timed out, and the local PubMed response was an interstitial rather than scientific content. medical\_source\_validation.json explicitly marks that response invalid for evidence; actual clinical text was read through official web retrieval. File existence was not mistaken for successful content acquisition.\par
All downloaded source bytes, author versions, repository files, replay scripts and images are confined to sources/0044 or verification/0044. provenance.json hashes the actual artifacts. The report PDF was rendered using the existing Lumen HTML/PDF engine from the same structured content as the self-contained TeX, then visually inspected after Poppler rendering. No local TeX compilation is claimed; controller compilation and final acceptance remain outstanding.\par
Primary AI work covered main-paper review, static source inspection, full-response/context checking, primary medical-source verification, integration and report/visual QA. Separate AI checks covered quantitative reconstruction, repository/version provenance and selected clinical/pragmatic interpretation. No human-review declaration was supplied on the user’s behalf. No author code, paid model calls, installations, external messages or controller edits occurred. Batch token:[internal identifier omitted].\par
\end{document}
